  \subsubsection{床関数を整数の個数として数える}\label{節：ガウス型の個数と格子点}
    \,$X\ge0$に対して$\lfloor X\rfloor$は,\,$1\le k\le X$を満たす整数$k$の個数である.
    例えば$\lfloor\frac nm\rfloor$は,\,$1,\ldots,n$に含まれる$m$の倍数の個数を表す.
    \bookterm{floor-function}{床関数}を含む和は,\textbf{何を数えているか}から意味を読み直せる.

    座標がともに整数である点を格子点という.
    例えば$\sum_{k=1}^{n}\lfloor ck\rfloor\ (c>0)$は,
    \,$x=k$を固定するたびに,\,$1\le y\le ck$を満たす整数$y$を数えている.
    同じ値の群で整理する方法とは別に,図形の中の点を数える方法としても理解できる.
    二つの使い方は重なり,前の群\bookterm{sequence}{数列}の別解も整数の組を数え直す方法であった.

    \noindent \bookblocklink{example-gauss-role-10}{problem}{answer}{【例題】}\par\nobreak
    \begin{enumerate}[label=(\arabic*),parsep=3pt,beginpenalty=10000,format=\bookitemlink{example-gauss-role-10}{problem}{answer}]
      \item \bookterm{sequence}{数列}を次のように定める.
            \[
              a_0=0,\qquad
              a_n=a_{n-1}+\left\lfloor\frac{5n}{2}\right\rfloor
              \quad(n\ge1).
            \]
            \bookterm{general}{一般項}を求め,\,$a_n$をある範囲の格子点の個数として説明せよ.
    \end{enumerate}

    \noindent \bookblocklink{example-gauss-role-10}{hint}{problem}{【方針】}\par\nobreak
    \begin{enumerate}[label=(\arabic*),parsep=3pt,beginpenalty=10000,format=\bookitemlink{example-gauss-role-10}{hint}{problem}]
      \item \,$\lfloor\frac{5k}{2}\rfloor=2k+\lfloor\frac k2\rfloor$と分けます.
            後者の和を偶奇で求め,\bookterm{floor-function}{床関数}が数えている整数$y$の範囲も確認しましょう.
    \end{enumerate}

    \noindent \bookblocklink{example-gauss-role-10}{answer}{problem}{【解答】}\par\nobreak
    \begin{enumerate}[label=(\arabic*),parsep=3pt,beginpenalty=10000,format=\bookitemlink{example-gauss-role-10}{answer}{problem}]
      \item \[
              a_n=n(n+1)+\left\lfloor\frac{n^2}{4}\right\rfloor.
            \]
            これは整数$x,y$で,
            \[
              1\le x\le n,\qquad 1\le y\le\frac{5x}{2}
            \]
            を満たす格子点の個数である.
    \end{enumerate}

    \noindent \bookblocklink{example-gauss-role-10}{solution}{problem}{【解法】}\par\nobreak
    \begin{enumerate}[label=(\arabic*),parsep=3pt,beginpenalty=10000,format=\bookitemlink{example-gauss-role-10}{solution}{problem}]
      \item \bookterm{difference}{階差}を足し,整数$2k$を\bookterm{floor-function}{床関数}の外に出すと,
            \begin{align*}
              a_n&=\sum_{k=1}^{n}\left\lfloor\frac{5k}{2}\right\rfloor\\
                 &=n(n+1)+\sum_{k=1}^{n}\left\lfloor\frac k2\right\rfloor.
            \end{align*}
            後ろの和は$n=2m$なら,
            \[
              0+1+1+2+2+\cdots+(m-1)+(m-1)+m=m^2,
            \]
            \,$n=2m+1$なら$m^2+m$である.
            従っていずれの場合も,
            \[
              \sum_{k=1}^{n}\left\lfloor\frac k2\right\rfloor
              =\left\lfloor\frac{n^2}{4}\right\rfloor
            \]
            となり,\bookterm{general}{一般項}を得る.\,$n=0$でも成立する.

            格子点としての意味も確かめよう.\,$x=k$と固定したとき,
            \,$1\le y\le\frac{5k}{2}$を満たす整数$y$は$\lfloor\frac{5k}{2}\rfloor$個である.
            \,$k=1,\ldots,n$について足した値が$a_n$だから,【解答】の範囲の点を数えている.
            \bookterm{sequence}{数列}の値,\bookterm{floor-function}{床関数}の和,図形の中の個数が同じ量を表している.
    \end{enumerate}
